\begin{conclusion}[Ellipsization of spectral quartic curves and isomorphism identification of minimal models]\label{con:xi-terminal-zm-elliptic-normalization}
Follow the conclusion \ref{con:xi-terminal-zm-order4-rational} mark, and order
\[
\mathcal C:=\bigl\{(\lambda,y)\in\mathbb A^2:\ \Pi(\lambda,y)=0\bigr\}.
\]
Define new variables
\[
w:=2y+1-2\lambda^{2}.
\]
And order
\[
u:=y-\lambda^{2}.
\]
then in \((x,u)\) coordinates (where \(x=\lambda\)), there is an identity under
\[
\Pi\bigl(x,\ u+x^{2}\bigr)=u^{2}+u-(x^{3}-x),
\]
thus in \(\QQ\) to obtain the minimum model with integer coefficients
\[
E_{\min}:\quad u^{2}+u=x^{3}-x,
\]
and \((x,u)\leftrightarrow(\lambda,y)\) is a reciprocal polynomial substitution, thus giving at the affine level \(\mathcal C\) and \(E_{\min}\) of \(\QQ\)-isomorphism;
If we take the projective closure, then \(\Pi(\lambda,y)=0\) Singular points may appear at infinity, but after normalization they become smooth elliptic curves.
This isomorphism can be replaced explicitly by
\[
(x,u)\longmapsto(\lambda,y)=(x,\ u+x^{2}),
\qquad
(\lambda,y)\longmapsto(x,u)=(\lambda,\ y-\lambda^{2})
\]
given directly.
further introduction \(w=2u+1\) (equivalent to the definition of the above formula \(w=2y+1-2\lambda^{2}\)), we get the isomorphism
\[
\mathcal C\ \dashrightarrow\ \mathcal E,\qquad
(\lambda,y)\longmapsto(\lambda,w),
\qquad
\mathcal E:\ w^2=4\lambda^{3}-4\lambda+1,
\]
Its inverse transformation is given by
\[
y=\frac{2\lambda^{2}+w-1}{2}
\]
given.
make \(X=\lambda,\ Y=w/2\),but \(\mathcal E\) is the short Weierstrass curve
\[
E:\ Y^2=X^3-X+\frac14,
\]
And in plane coordinates \((\lambda,y)\) exists a completely explicit decomposition of squared differences on
\[
\Pi(\lambda,y)=\Bigl(y-\lambda^2+\tfrac12\Bigr)^2-\Bigl(\lambda^3-\lambda+\tfrac14\Bigr).
\]
Therefore let \(Y:=y-\lambda^2+\tfrac12\)(that is \(Y=w/2\)) can obtain the same function domain
\[
\QQ(\mathcal C)\cong \QQ(E),\qquad \lambda=X,\quad Y=y-\lambda^2+\tfrac12,
\]
and gives a fully explicit birational reciprocal inversion on affine smooth open sets
\[
\nu:\ E\to \mathcal C,\quad (X,Y)\longmapsto(\lambda,y)=(X,\ X^2+Y-\tfrac12),
\qquad
\nu^{-1}:\ \mathcal C\dashrightarrow E,\quad (\lambda,y)\longmapsto(X,Y)=\bigl(\lambda,\ y-\lambda^2+\tfrac12\bigr).
\]
Its discriminant is the same as \(j\)-invariants are respectively
\[
\Delta(E)=37,\qquad j(E)=\frac{48^{3}}{37}=\frac{110592}{37}.
\]
further order \(y_{\mathrm{int}}:=Y-\tfrac12\)(that is \(y_{\mathrm{int}}=u=y-\lambda^{2}\)), then the minimum model with integer coefficients is obtained
\[
E_{\min}:\quad y_{\mathrm{int}}^{2}+y_{\mathrm{int}}=X^{3}-X,
\]
That \(\Delta(E_{\min})=37\) and \(j(E_{\min})=48^{3}/37\)。
therefore \((\lambda,y)\)-The algebraic branch structure of spectral curves can be equivalently promoted to elliptic curves \(E\), thus converting the bifurcation point, analytic continuation and singularization problems into \(E\) on algebraic geometric data.
In particular,\(\Delta(E)=37\) indicates that the elliptic curve divides \(37\) is well reduced everywhere; and corresponds to homologous classes under the Cremona/LMFDB flag \(37\mathrm a\) of the strong Weil curve \(37\mathrm a1\) (derivative \(37\), see notes \ref{rem:xi-terminal-elliptic-37a1-modular-anchor}）。
\par
Under the above isomorphism, the weight parameter is as \(E\) can be written as
\[
\boxed{\;
y=\frac{2\lambda^{2}+w-1}{2}=X^2+Y-\frac12\ \in\ \QQ(E)\; } ,
\qquad (O\text{ for }E\text{ the infinity point}).
\]
but \(y:E\to\mathbb P^1_y\) is the number of times \(4\), and the divisor of zero fibers and unit fibers is completely explicit:
\begin{enumerate}
  \item consists of \(y=0\iff w=1-2\lambda^{2}\) have to
  \[
  (1-2\lambda^2)^2=4\lambda^3-4\lambda+1
  \ \Longleftrightarrow\
  4\lambda(\lambda-1)^2(\lambda+1)=0,
  \]
  thereby
  \[
  \boxed{\;
  \mathrm{div}(y)=\Bigl(0,\frac12\Bigr)+2\Bigl(1,-\frac12\Bigr)+\Bigl(-1,-\frac12\Bigr)-4\cdot O\; }.
  \]
  \item consists of \(y=1\iff w=3-2\lambda^{2}\) have to
  \[
  (3-2\lambda^2)^2=4\lambda^3-4\lambda+1
  \ \Longleftrightarrow\
  4(\lambda-2)(\lambda-1)(\lambda+1)^2=0,
  \]
  thereby
  \[
  \boxed{\;
  \mathrm{div}(y-1)=\Bigl(2,-\frac52\Bigr)+\Bigl(1,\frac12\Bigr)+2\Bigl(-1,\frac12\Bigr)-4\cdot O\; }.
  \]
\end{enumerate}
By the main divisor in \(\mathrm{Pic}^0(E)\cong E\) corresponds to the unit element, which can also be equivalently written as two group law relations
\[
2\Bigl(1,-\frac12\Bigr)+\Bigl(-1,-\frac12\Bigr)+\Bigl(0,\frac12\Bigr)=O,
\qquad
2\Bigl(-1,\frac12\Bigr)+\Bigl(1,\frac12\Bigr)+\Bigl(2,-\frac52\Bigr)=O.
\]
Combined with conclusion \ref{con:xi-terminal-zm-elliptic-mw-generator-alignment} generator of \(R=(0,\tfrac12)\) aligned with the low magnification point, the decomposition of the above two fibers can also be written in the sense of divisor as
\[
y^{-1}(0)=\{R\}+2\{-2R\}+\{3R\},\qquad
y^{-1}(1)=\{2R\}+2\{-3R\}+\{4R\},
\]
Each point is counted by weight.
therefore \(y\) exist \(y=0\) and \(y=1\) appears one double fiber point each (respectively located at \(\lambda=1\) and \(\lambda=-1\)), while the remaining divergence values ​​are decomposed by the discriminant (Conclusion \ref{con:xi-terminal-zm-discriminant-leyang} Or equivalent result expression characterization, see conclusion \ref{con:xi-terminal-zm-leyang-cubic-branch-image}); this gives a "divisor-divergence" certificate of pure algebraic geometry.
\end{conclusion}
\begin{proof}[Points of proof]
Will \(y=(2\lambda^2+w-1)/2\) Substitute \(\Pi(\lambda,y)\) and simplify to get the identity
\[
4\,\Pi(\lambda,y)=w^2-\bigl(4\lambda^3-4\lambda+1\bigr),
\]
This leads to the birational equivalence shown.
Take the Weierstrass form \(a=-1,b=1/4\),but
\(
\Delta=-16(4a^3+27b^2)=37
\),and \(j=c_4^3/\Delta\)(in \(c_4=-48a\)) gives \(j(E)=110592/37\)。
Reproducible verification script:\texttt{scripts/exp\_fold\_zm\_elliptic\_leyang\_arithmetic\_audit.py}. Certification completed.
\end{proof}

% (User input window)
% (Source tag) /killo-golden
% Timestamp: 2026-02-18 15:11:15 (Asia/Singapore)
%
% Content is rewritten using the existing notation and auditable interfaces of this paper,
% without introducing manual numbering.

\begin{conclusion}[Unique singularity of the spectral curve at infinity: an \texorpdfstring{$A_4$}{A4}-type unibranch cusp with \texorpdfstring{$\delta=2$}{delta=2}]\label{con:xi-terminal-zm-spectral-curve-infty-a4-cusp}
Retain the spectral quartic from Conclusion \ref{con:xi-terminal-zm-elliptic-normalization}:
\[
\Pi(\lambda,y)=\lambda^4-\lambda^3-(2y+1)\lambda^2+\lambda+y(y+1)\in\QQ[\lambda,y],
\]
and let \(\overline{\mathcal C}\subset \mathbb P^1_\lambda\times \mathbb P^1_y\) denote the closure of the affine plane curve \(\Pi(\lambda,y)=0\) (viewed as a curve of type \((4,2)\)).
Then the following hold:
\begin{enumerate}
  \item \(\overline{\mathcal C}\) intersects the boundary at infinity only at \((\lambda,y)=(\infty,\infty)\), and this point is the unique singular point of \(\overline{\mathcal C}\).
  \item Taking local coordinates near \((\infty,\infty)\),
  \[
  \mu:=\frac{1}{\lambda},\qquad v:=\frac{1}{y},
  \]
  the local equation of \(\overline{\mathcal C}\) can be written as a quadratic in \(v\), and its discriminant has the strict factorization
  \[
  \boxed{\;\Delta_v(\mu)=\mu^5\cdot(\mu^3-4\mu^2+4)\;}
  \]
  so the singularity is analytically equivalent to the standard form
  \[
  v^2=u^5,
  \]
  i.e. an \(\,A_4\) (ramphoid cusp) unibranch cusp; hence its \(\delta\)-invariant is
  \[
  \boxed{\;\delta=\frac{(2-1)(5-1)}{2}=2\; }.
  \]
  \item A smooth \((4,2)\)-curve has arithmetic genus \((4-1)(2-1)=3\); since the unique singularity above contributes \(\delta=2\), the geometric genus of the normalization is \(3-2=1\), consistent with the identification ``normalization is an elliptic curve'' in Conclusion \ref{con:xi-terminal-zm-elliptic-normalization}.
\end{enumerate}
This \(A_4\) cusp can also be viewed geometrically as a planar compressed encoding of the quadruple totally ramified pole structure at \(y=\infty\) described in Conclusion \ref{con:xi-terminal-zm-elliptic-y-hurwitz-passport}.
\end{conclusion}
\begin{proof}[Proof sketch]
Conclusion \ref{con:xi-terminal-zm-elliptic-normalization} already gives that the affine part of \(\Pi(\lambda,y)=0\) is isomorphic to a smooth elliptic curve,
so there are no singularities at finite points; hence any singularity must lie on the boundary at infinity.
Viewing \(\Pi\) as a bihomogeneous equation on \(\mathbb P^1_\lambda\times \mathbb P^1_y\), one checks directly that neither \(\lambda=\infty\), \(y\neq\infty\) nor \(y=\infty\), \(\lambda\neq\infty\) can satisfy the equation,
thus the boundary intersection point must be \((\infty,\infty)\).
\par
Near \((\infty,\infty)\), set \(\mu=1/\lambda\), \(v=1/y\). Then
\[
\Pi\!\left(\frac{1}{\mu},\frac{1}{v}\right)=0
\quad\Longleftrightarrow\quad
F(\mu,v)=0,
\]
where multiplying by \(\mu^4v^2\) clears denominators and yields
\[
F(\mu,v)
=(1-\mu-\mu^2+\mu^3)v^2+\mu^2(\mu^2-2)v+\mu^4.
\]
Treating \(F\) as a quadratic in \(v\), its discriminant is
\[
\Delta_v(\mu)=\mu^4(\mu^2-2)^2-4\mu^4(1-\mu-\mu^2+\mu^3)
=\mu^5(\mu^3-4\mu^2+4),
\]
giving the displayed factorization.
Completing the square gives
\[
\bigl(2(1-\mu-\mu^2+\mu^3)v+\mu^2(\mu^2-2)\bigr)^2=\Delta_v(\mu).
\]
Since \(\mu^3-4\mu^2+4\) is a unit at \(\mu=0\) (value \(4\)), its unit factor can be absorbed by analytic coordinate change,
hence the singularity is analytically equivalent to \(v^2=u^5\); \(\delta\) follows from the standard formula for the plane branch \(v^2=u^5\).
Finally, the genus calculation is the arithmetic genus minus \(\delta\). QED.
\end{proof}




\begin{conclusion}[frequency \(4\) linear system \texorpdfstring{$|4O|$}{|4O|} Projective normal model of: Explicit intersection of two quadratic surfaces]\label{con:xi-terminal-zm-elliptic-degree4-projective-model}
Follow the conclusion \ref{con:xi-terminal-zm-elliptic-normalization} The short Weierstrass model
\(
E:\ Y^{2}=X^{3}-X+\tfrac14
\),remember \(O\) is an infinite point, and let
\(
y:=X^2+Y-\tfrac12\in\QQ(E)
\)。
but \(y\) exist \(O\) has a fourth-order pole, so \(1,X,X^{2},y\in H^{0}\!\bigl(E,\mathcal O_{E}(4O)\bigr)\)；
and they constitute \(H^{0}(E,\mathcal O_{E}(4O))\) is a set of basis.
By a completely linear system \(|4O|\) induced morphism
\[
\kappa:E\hookrightarrow \mathbb P^{3},\qquad
P\longmapsto [U_{0}:U_{1}:U_{2}:U_{3}]=[1:X(P):X(P)^{2}:y(P)]
\]
is a projective embedding, and its image is the degree \(4\) is a projective normal curve, and in \(\mathbb P^{3}\) is described by two quadratic equations:
\[
U_{1}^{2}-U_{0}U_{2}=0,
\qquad
(U_{2}-U_{3})^{2}-U_{0}(U_{2}-U_{3})-U_{1}U_{2}+U_{1}U_{0}=0.
\]
In this model, the projection \([U_{0}:U_{3}]\) and rational functions \(y\) the number of times defined \(4\) overlay \(E\to\mathbb P^{1}_{y}\) consistent.
\end{conclusion}
\begin{proof}[Points of proof]
By conclusion \ref{con:xi-terminal-zm-elliptic-y-hurwitz-passport}，\(y\) exist \(O\) there is only one fourth-order pole, so
\(
y\in H^{0}(E,\mathcal O_{E}(4O))
\)。
And because of \(\mathrm{ord}_{O}(X)=-2\)、\(\mathrm{ord}_{O}(Y)=-3\), it can be seen \(1,X,X^{2}\in H^{0}(E,\mathcal O_{E}(4O))\),and
\(
y-X^{2}=Y-\tfrac12
\)
exist \(O\), the extreme order is \(3\)。
like \(a_{0}+a_{1}X+a_{2}X^{2}+a_{3}y=0\),Compare \(O\) at the extreme level comes first served \(a_{3}=-a_{2}\), and then by \(y-X^{2}=Y-\tfrac12\)'s third-order pole derivation \(a_{2}=a_{3}=0\), then \(a_{1}=a_{0}=0\), so the four are linearly independent; since \(\dim H^{0}(E,\mathcal O_{E}(4O))=4\), they form a set of bases.
when \(\deg(4O)=4\ge 3\) hour \(|4O|\) is very ample, so \(\kappa\) is an embedding; its image is an elliptic normal curve, so it is the complete intersection of two quadratic surfaces.
The two quadratic equations can be checked directly by algebraic elimination: in the affine graph \(U_{0}=1\), the first equation is the identity \(X^{2}=X^{2}\)；
The second formula becomes
\[
(X^{2}-y)^{2}-(X^{2}-y)-X^{3}+X=0,
\]
Will \(y=X^{2}+Y-\tfrac12\) and use the curve equation \(Y^{2}=X^{3}-X+\tfrac14\) can be obtained by simplifying.
The last projection statement is defined by \(U_{3}/U_{0}=y\) is launched directly. Certification completed.
\end{proof}

\begin{conclusion}[Primary divisor of conjugate weight branch: complete writing method on low-magnification rational point orbit]\label{con:xi-terminal-zm-elliptic-conjugate-weight-principal-divisors}
Follow the conclusion \ref{con:xi-terminal-zm-elliptic-normalization} notation, and let the ellipse be inverted as
\(
\iota:(X,Y)\mapsto(X,-Y)
\)。
Define the conjugate weight function
\[
y^{\iota}:=\iota(y)=X^{2}-Y-\frac12\in\QQ(E).
\]
remember \(R:=(0,\tfrac12)\in E(\QQ)\) is the conclusion \ref{con:xi-terminal-zm-elliptic-mw-generator-alignment} The corresponding point of the free generator.
but \(y^{\iota}\) and \(y^{\iota}-1\) satisfies the principal divisor of
\[
\operatorname{div}(y^{\iota})=(-R)+(-3R)+2(2R)-4O,
\qquad
\operatorname{div}(y^{\iota}-1)=(-4R)+(-2R)+2(3R)-4O.
\]
Therefore two special fibers (\(y^{\iota}=0\) and \(y^{\iota}=1\)) zero points (including multiplicity) are all supported by \(R\) on the generated Mordell--Weil orbit;
and \(y^{\iota}\)(respectively \(y^{\iota}-1\)) as \(|4O|\), the distribution of its zero divisors on this orbit is uniquely fixed by the above-mentioned principal divisor identity.
\end{conclusion}
\begin{proof}[Points of proof]
Depend on \(y^{\iota}=X^{2}-Y-\tfrac12\) and conclusion \ref{con:xi-terminal-zm-elliptic-normalization} of the subcertificate,
right \(\operatorname{div}(y)\)、\(\operatorname{div}(y-1)\) impose \(\iota\)(Notice \(\iota(O)=O\)) can be obtained
\[
\operatorname{div}(y^{\iota})=\Bigl(0,-\tfrac12\Bigr)+\Bigl(-1,\tfrac12\Bigr)+2\Bigl(1,\tfrac12\Bigr)-4O,
\qquad
\operatorname{div}(y^{\iota}-1)=\Bigl(2,\tfrac52\Bigr)+\Bigl(1,-\tfrac12\Bigr)+2\Bigl(-1,-\tfrac12\Bigr)-4O.
\]
And then the conclusion \ref{con:xi-terminal-zm-elliptic-mw-generator-alignment} Alignment of low magnification points
\(
[2]R=(1,\tfrac12),\ [3]R=(-1,-\tfrac12),\ [4]R=(2,-\tfrac52)
\)
The above coordinate representation can be rewritten as shown \(\langle R\rangle\) orbital form. Certification completed.
\end{proof}

\begin{conclusion}[Necessary and sufficient conditions for rational spectrum roots:\texorpdfstring{$\Pi(\lambda,y)$}{Pi(lambda,y)} exist \texorpdfstring{$\mathbb Q$}{Q} There is a foundation if and only if \(y\in y\bigl(E(\QQ)\bigr)\)]\label{con:xi-terminal-zm-pi-rational-root-iff-y-image}
Follow the conclusion \ref{con:xi-terminal-zm-elliptic-normalization} spectrum four times
\(
\Pi(\lambda,y)=0
\)
and its elliptical isomorphism.
Then for any \(y_{0}\in\QQ\), the following two conditions are equivalent:
\begin{enumerate}
  \item exists \(\lambda\in\QQ\) makes \(\Pi(\lambda,y_{0})=0\)；
  \item exists \(P\in E(\QQ)\) makes \(y_{0}=y(P)\)。
\end{enumerate}
And when the above conditions are met, it is advisable \(\lambda=X(P)\) as the corresponding rational spectral root.
In particular, at this time \(\Pi(\lambda,y_{0})\) exist \(\QQ[\lambda]\) contains rational linear factors \(\lambda-X(P)\)；
If you will \(X(P)\) written as a reduced fraction \(X(P)=U/V^{2}\)（\(U,V\in\ZZ,\ \gcd(U,V)=1,\ V>0\)), then there exists a cubic polynomial \(G(\lambda)\in\QQ[\lambda]\) makes
\[
\Pi(\lambda,y_{0})=(V^{2}\lambda-U)\,G(\lambda),
\]
And the denominator satisfies
\[
\mathrm{den}(y_{0})=\mathrm{den}\bigl(X(P)\bigr)^{2}=V^{4}.
\]
\end{conclusion}
\begin{proof}[Points of proof]
if exists \(\lambda\in\QQ\) make \(\Pi(\lambda,y_{0})=0\), from the conclusion \ref{con:xi-terminal-zm-elliptic-normalization} The identity of
\(
4\Pi(\lambda,y)=w^{2}-(4\lambda^{3}-4\lambda+1)
\)
(in \(w=2y+1-2\lambda^{2}\)) It can be known \((\lambda,w)\in\mathcal E(\QQ)\),thereby \(P=(X,Y)=(\lambda,w/2)\in E(\QQ)\), and satisfy \(y(P)=y_{0}\)。
On the contrary, if \(P\in E(\QQ)\),but \((\lambda,y)=(X(P),y(P))\) satisfy \(\Pi(\lambda,y)=0\),thereby \(\lambda=X(P)\) gives \(\Pi(\lambda,y_{0})\) of rational roots.
\par
Linear Factors and Clear Denominator Factors \((V^{2}\lambda-U)\) Depend on \(\lambda=X(P)=U/V^{2}\) is launched immediately for root.
The denominator square relationship can be directly expressed by the curve equation and \(y=X^{2}+Y-\tfrac12\) launch: write \(X=U/V^{2}\) (contracted), then by
\(
Y^{2}=X^{3}-X+\tfrac14
\)
have to \(Y=W/(2V^{3})\)（\(W\in\ZZ\))and \(W^{2}=4U^{3}-4UV^{4}+V^{6}\)；
thereby
\[
y(P)=\frac{U^{2}}{V^{4}}+\frac{W}{2V^{3}}-\frac12=\frac{2U^{2}+WV-V^{4}}{2V^{4}}.
\]
Notice \(2U^{2}+WV-V^{4}\) is always an even number, so \(y(P)=N/V^{4}\)（\(N\in\ZZ\)）。
odd prime numbers \(p\mid V\),but \(WV\equiv V^{4}\equiv 0\pmod p\) and \(2N\equiv 2U^{2}\not\equiv 0\pmod p\),thereby \(p\nmid N\)；
like \(2\mid V\),but \(U\) is an odd number;
and by \(W^{2}=4U^{3}-4UV^{4}+V^{6}\) in mold \(8\) \(W^{2}\equiv 4\pmod 8\), so \(W=2W_{1}\) and \(W_{1}\) is an odd number.
then \(WV\equiv 0\pmod 4\)、\(V^{4}\equiv 0\pmod 4\),and \(2U^{2}\equiv 2\pmod 4\),thereby \(2N\equiv 2\pmod 4\),Right now \(N\) is an odd number;
To sum up \(\gcd(N,V)=1\)。
therefore \(\mathrm{den}(y(P))=V^{4}=\mathrm{den}(X(P))^{2}\). Certification completed.
\end{proof}

\begin{remark}[Double discriminant bridge:\texorpdfstring{$\Pi$}{Pi} two discriminant directions]\label{rem:xi-terminal-zm-pi-double-discriminant-bridge}
Will \(\Pi(\lambda,y)\) is deemed to be about \(y\) is a quadratic polynomial, then its \(y\)-the discriminant is
\[
\mathrm{Disc}_{y}\bigl(\Pi(\lambda,y)\bigr)=4\lambda^{3}-4\lambda+1.
\]
Equivalently, let \(w:=2y+1-2\lambda^{2}\),but
\[
4\Pi(\lambda,y)=w^{2}-\bigl(4\lambda^{3}-4\lambda+1\bigr).
\]
So the same equation is in \(y\)-direction discriminant gives the elliptic curve model \(\mathcal E:w^{2}=4\lambda^{3}-4\lambda+1\) (see conclusion \ref{con:xi-terminal-zm-elliptic-normalization}）。
On the other hand, the \(\Pi(\lambda,y)\) is deemed to be about \(\lambda\) of a fourth degree polynomial, then its \(\lambda\)-the discriminant is in \(\ZZ[y]\) is decomposed into
\[
\mathrm{Disc}_{\lambda}\bigl(\Pi(\lambda,y)\bigr)=-y(y-1)P_{\mathrm{LY}}(y),
\qquad
P_{\mathrm{LY}}(y):=256y^{3}+411y^{2}+165y+32,
\]
thereby obtaining the attribute \(2\)’s discriminant curve \(C:\ W^{2}=-y(y-1)P_{\mathrm{LY}}(y)\)(note here \(W\) to avoid conflicts with elliptical coordinates \(w\) confusion).
This "quadratic-quadratic" dual-discriminant structure uniformly encodes the ellipse normalization, bifurcation images, and the quadratic resolution of the quadratic split domain in the same \(\Pi\) among.
\end{remark}

\begin{remark}[External naming and cross-validation interface:\texorpdfstring{$37\mathrm a1$}{37a1}、\texorpdfstring{$X_0^+(37)$}{X0+(37)} with level \(37\) new form]\label{rem:xi-terminal-elliptic-37a1-modular-anchor}
The elliptic curve of this segment \(E: Y^2=X^3-X+\tfrac14\) exist \(\QQ\) and the integral coefficient model
\[
E_0:\quad \eta^2+\eta=x^3-x
\]
isomorphism; explicit isomorphism is given by
\(
X=x,\ Y=\eta+\tfrac12
\)
given.
This article uniformly adopts the standard convention of Weierstrass invariants.
\[
c_4^3-c_6^2=1728\Delta,
\]
and in the general Weierstrass form
\(
\eta^2+a_1x\eta+a_3\eta=x^3+a_2x^2+a_4x+a_6
\)
The invariant notation of \(E_0\) Pick
\(
(a_1,a_2,a_3,a_4,a_6)=(0,0,1,-1,0)
\)，
thereby \(b_2=0\)、\(b_4=-2\)、\(b_6=1\)、\(b_8=-1\), and obtained by the standard discriminant formula \(\Delta(E_0)=64-27=37\), and the short Weierstrass model above \(E\) are calculated consistently.
If a discriminant token with the opposite sign appears in external literature, it corresponds to the \(\Delta\) adopts different symbol conventions and does not change the purpose of this article. \(\Delta(E)=37\) is the anchoring follow-up argument.
according to \cite{LMFDBEllipticCurve37a1} records (accessed at \(2026\text{-}02\text{-}17\)), under the standard identity of Cremona/LMFDB,\(E\) is the derivation \(37\) homologous class \(37\mathrm a\) of the strong Weil curve (notation \(37\mathrm a1\)）；
its minimal integer coefficient model \(E_0:\eta^2+\eta=x^3-x\) can also be found with Cremona in \(N=37\).\cite{CremonaExamplesPDF}
\par
The above external identifier does not participate in the subsequent derivation of this article, but provides a comparison interface with a classic background: the weight corresponding to the homologous class \(2\)、level \(37\) The new form orbit is \(37.2.a.a\), and for any good element \(p\neq 37\) have
\(
a_p=p+1-\#E(\FF_p)
\)
is consistent with the Hecke eigenvalues ​​of this new form, so that \(L(E,s)=L(f,s)\) (modularity theorem; can also be understood in Langlands caliber as \(\mathrm{GL}_2(\mathbb{A}_\QQ)\) corresponding automorphic representation of docking).\cite{LMFDBNewform372aa}
\par
So for every prime number \(\ell\), elliptic curve \(E\) of \(\ell\)-advanced Tate module induced two-dimensional Galois representation
\[
\rho_{E,\ell}:\mathrm{Gal}(\overline{\QQ}/\QQ)\longrightarrow \mathrm{GL}_2(\ZZ_\ell),
\]
its only in \(\{37,\ell\}\) may differ; and for any good \(p\neq 37,\ell\), the characteristic polynomial of the Frobenius element is \(T^2-a_pT+p\), giving local Euler factors consistent with the new form.
\par
And from Sturm boundedness, it can be seen that the external anchor has finite certificate: \(k=2,N=37\) have \([SL_2(\ZZ):\Gamma_0(37)]=38\), so Sturm bound
\(
B=\bigl\lfloor \frac{k}{12}[SL_2(\ZZ):\Gamma_0(N)]\bigr\rfloor
\)
equal \(6\); thus only need to verify up to \(q^6\) can uniquely determine this new form of orbital.\cite{LMFDBSturmBound}
Under the mold curve diameter,\(X_0(37)\) is a genus \(2\) curve; LMFDB record \(E_0\) can be used as a quotient of its Fricke involution
\(
X_0^+(37):=X_0(37)/\langle w_{37}\rangle
\)
One of the algebraic models.\cite{LMFDBEllipticCurve37a1}
\par
make \(q=e^{2\pi i\tau}\)。
Under the above external anchoring, a new form of normalization is desirable \(f(\tau)=\sum_{n\ge 1}a_n q^n\) makes
\(
f(q)\,\dd q/q
\)
equal \(E_0\) of N\Pullback of 'eron differential under modular parameterization;
So the weight function \(y=\eta+x^2\in\QQ(E_0)\) as \(X_0^+(37)\) has a principal part at the infinity cusp \(q^{-4}\), its truncated expansion corresponds to the new form \(q\)-expanded to see the following formula:
\input{sections/generated/eq_fold_zm_elliptic_modular_y_qexp_audit}
Reproducible verification script:\texttt{scripts/exp\_fold\_zm\_elliptic\_modular\_y\_qexp\_audit.py}。
\end{remark}

\begin{conclusion}[\texorpdfstring{$j$}{j}-Standard basic domain inversion of invariants: pure virtual module parameters \texorpdfstring{$\tau$}{tau} Uniqueness and Numerical Certificate]\label{con:xi-terminal-zm-elliptic-j-inversion-tau}
in conclusion \ref{con:xi-terminal-zm-elliptic-normalization} given
\(
j(E)=110592/37
\)。
then in \(SL_{2}(\ZZ)\) within the standard basic domain, satisfying \(j(\tau)=j(E)\) modular parameters \(\tau\) is unique and purely imaginary
\[
\tau=i\cdot 1.221127360764627\ldots.
\]
Correspondingly
\(
q
\)-parameter is
\[
q=e^{2\pi i\tau}=0.000465420392294374\ldots.
\]
therefore \(E(\CC)\) can be explicitly homogenized to \(\CC/(\ZZ+\tau\ZZ)\); and the weight function \(y\) under this homogenization is the order \(4\), its analytic scale and discriminant prime \(37\) are consistent with the modular anchor points.
\end{conclusion}
\begin{proof}[Points of proof]
because \(j(E)\in\RR\) and \(j(E)>1728=j(i)\), standard module theory gives that the only solution in the basic domain must fall on the pure imaginary axis, that is \(\tau=it\)（\(t>1\)）。
pair function \(t\mapsto j(it)\) can be recovered by performing numerical inversion \(t\) and \(q=e^{-2\pi t}\)。
The inversion can be performed with high accuracy through the Eisenstein series \(q\)-expand implementation, and by script \texttt{scripts/exp\_fold\_zm\_elliptic\_modular\_y\_qexp\_audit.py} Recurrence verification. Certification completed.
\end{proof}

\begin{conclusion}[Identity for the number of modular fiber roots: given by \texorpdfstring{$\Pi(\lambda,y)$}{Pi(lambda,y)} Decomposed statistical recovery of \texorpdfstring{$a_p$}{a\_p}]\label{con:xi-terminal-zm-fiber-root-counting-hecke}
Take prime number \(p\notin\{2,37\}\), and in \(\FF_p\) consider affine curves on
\[
\mathcal C_p:=\{(\lambda,y)\in\FF_p^2:\ \Pi(\lambda,y)=0\}.
\]
for each \(y\in\FF_p\) defines the number of fibers
\[
N_p(y):=\card{\{\lambda\in\FF_p:\ \Pi(\lambda,y)=0\}}.
\]
Then there is a strict point counting identity
\[
\boxed{\;
\card{E(\FF_p)}=1+\sum_{y\in\FF_p}N_p(y)\;}
\]
The one on the right \(1\) corresponds to the point at infinity \(O\)。
make \(a_p:=p+1-\card{E(\FF_p)}\) is the Frobenius trace (see note \ref{rem:xi-terminal-elliptic-37a1-modular-anchor}), then equivalently
\[
\boxed{\;
a_p=p-\sum_{y\in\FF_p}N_p(y)\;}
\]
And get an effective estimate from the Hasse bound
\[
\left|\frac{1}{p}\sum_{y\in\FF_p}N_p(y)-1\right|
=\left|\frac{a_p}{p}\right|
\le \frac{2}{\sqrt p}.
\]
\end{conclusion}
\begin{proof}[Points of proof]
when \(p\notin\{2,37\}\), the conclusion \ref{con:xi-terminal-zm-elliptic-normalization} birational transformation of
\[
(\lambda,y)\longmapsto (\lambda,w),\qquad w:=2y+1-2\lambda^{2}
\]
exist \(\FF_p\) is still reversible (because \(2\) reversible), and put \(\mathcal C_p\) bijection to affine elliptic curve
\(
\mathcal E_p:\ w^2=4\lambda^3-4\lambda+1
\)
of \(\FF_p\)-point set; the affine model and the projective elliptic curve \(E\) differs only at infinity \(O\)。
therefore
\[
\sum_{y\in\FF_p}N_p(y)=\card{\mathcal C_p}=\card{E(\FF_p)}-1,
\]
This leads to the identity shown.
Substitute it into \(a_p:=p+1-\card{E(\FF_p)}\) \(a_p=p-\sum_{y\in\FF_p} N_p(y)\)。
The last expression is bounded by Hasse \(|a_p|\le 2\sqrt p\) is given.
Reproducible verification script:\texttt{scripts/exp\_fold\_zm\_elliptic\_fiber\_root\_counting\_audit.py}. Certification completed.
\end{proof}

% (User input window)

\begin{conclusion}[Decomposed equal distribution on tamed primes:\texorpdfstring{$S_4$}{S4} conjugate class density and square root error]\label{con:xi-terminal-zm-s4-modp-splitting-density}
Following the theorem \ref{thm:xi-terminal-zm-leyang-elliptic-structure} and conclusion \ref{con:xi-terminal-zm-generic-galois-s4}, the geometric single-valued group \(\mathrm{Mon}(\pi_y)\cong S_4\),
And order
\(
B:=\{0,1,\infty\}\cup\{P_{\mathrm{LY}}(y)=0\}\subset \mathbb P^1_y
\)
is the set of divergent values ​​covering \(\pi_y:E\to\mathbb P^1_y\) (see conclusion \ref{con:xi-terminal-zm-discriminant-leyang}).
Take the prime \(p\notin\{2,3,31,37\}\) and consider the undiverged parameter set on \(\FF_p\)
\[
U_p:=\FF_p\setminus\bigl(B\cap \FF_p\bigr)
=\Bigl\{y_0\in\FF_p:\ y_0(y_0-1)P_{\mathrm{LY}}(y_0)\neq 0\Bigr\}.
\]
Then when \(y_0\) is uniformly distributed on \(U_p\),
The decomposition form of the fourth-order polynomial \(\Pi(\lambda,y_0)\in\FF_p[\lambda]\) is equally distributed according to the conjugate class density of \(S_4\), and has a function domain Weil-type error of \(O(p^{-1/2})\);
More specifically, corresponding to the "decomposition type \(\leftrightarrow\) cyclic type of Frobenius on four roots", the limiting density of the five decomposition types is
\begin{itemize}
\item Irreducible to fourth degree (cyclic type \((4)\)): density \(6/24=1/4\);
\item One time \(\times\) cannot be reduced to three times (cyclic type \((3+1)\)): density \(8/24=1/3\);
\item Two irreducible quadratics (cyclic type \((2+2)\)): density \(3/24=1/8\);
\item Two linear \(\times\) One irreducible quadratic (cyclic type \((2+1+1)\)): Density \(6/24=1/4\);
\item Full split (cyclic type \((1+1+1+1)\)): density \(1/24\).
\end{itemize}
This conclusion is compatible with the mean fiber root number identity of the conclusion \ref{con:xi-terminal-zm-fiber-root-counting-hecke}: on undiverged fibers, the \(\FF_p\)-root number mean of \(\Pi(\lambda,y_0)\) tends to \(1\), with its global deviation controlled by \(a_p\).
\end{conclusion}
\begin{proof}[Proof sketch]
For \(p\notin\{2,3,31,37\}\), the conclusion \ref{con:xi-terminal-zm-discriminant-leyang} ensures that \(B\) still has six mutually different divergence values ​​(including \(\infty\)) under the module \(p\), and the divergence indexes are all relatively prime with \(p\), thus covering the tame divergence on \(U_p\).
From the conclusion that \(S_4\) is unique in \ref{con:xi-terminal-zm-generic-galois-s4}, take its \(S_4\)-Galois closure on \(\FF_p\) and apply the function domain Chebotarev theorem to the rational points of \(U_p\),
It can be obtained that the Frobenius conjugate classes are equally distributed in \(S_4\) according to the class size \(|C|/24\); the error term is given by the Weil bound of the corresponding Galois covering curve, and the magnitude is \(O(\sqrt p)\), which is \(O(p^{-1/2})\) when converted to a scale.
The correspondence between the cyclic form and the decomposition form of \(\FF_p[\lambda]\) is a standard fact.
Reproducible verification script: \texttt{scripts/exp\_fold\_zm\_elliptic\_fiber\_root\_counting\_audit.py} (the output includes decomposed frequency auditing). Certification completed.
\end{proof}

\begin{conclusion}[Precise control of collisions between covered minimal bad reduction prime sets and divergence values]\label{con:xi-terminal-zm-cover-minimal-bad-reduction-set}
Following the conclusion \ref{con:xi-terminal-zm-discriminant-leyang}'s divergence value set \(B\subset\mathbb P^1_y\), its integral coefficient model is given by
\[
y(y-1)P_{\mathrm{LY}}(y)=0
\]
is given (where \(P_{\mathrm{LY}}(y)\) is a cubic factor of \(\mathrm{Disc}_\lambda(\Pi)\)).
Then the set of prime numbers \(\{2,3,31,37\}\) constitutes the minimum bad reduction set covered by the spectrum, and the meaning is as follows:
\begin{itemize}
\item If \(p\notin\{2,3,31,37\}\), then \(B\) remains mutually different under the module \(p\), and covering \(\FF_p\) is tame and the Hurwitz passport \((2,2,2,2,2,4)\) remains unchanged, thus giving a good reduction;
\item If \(p=3\), then \(P_{\mathrm{LY}}(y)\equiv (y-1)^3\ (\mathrm{mod}\ 3)\), three non-trivial divergence values ​​collide in triple under the modulo \(3\), thus covering the divergence data degradation;
\item If \(p=31\), then
  \[
  P_{\mathrm{LY}}(y)\equiv 8\,(y-12)^{2}(y-6)\pmod{31},
  \]
As a result, double roots appear under module \(31\) and lead to double collisions of divergence values, covering the degradation of divergence data;
\item If \(p=37\), then
  \[
  P_{\mathrm{LY}}(y)\equiv -3\,(y-14)(y-6)^{2}\pmod{37},
  \]
And the elliptic curve \(E\) undergoes multiplicative bad reduction (\(\Delta(E)=37\)); therefore, there are both divergence value collisions and base curve bad reduction, and the coverage is degraded.
\end{itemize}
Therefore at all \(p\notin\{2,3,31,37\}\), the spectrum covers the "structurally stable" domestication case.
\end{conclusion}
\begin{proof}[Proof sketch]
From the conclusion \ref{con:xi-terminal-zm-discriminant-leyang},
\(
\mathrm{Disc}_\lambda(\Pi)=-y(y-1)P_{\mathrm{LY}}(y)
\)，
Therefore, the merging of divergent values ​​occurs if and only if the modulo \(p\) reduction of \(P_{\mathrm{LY}}\) is no longer separable.
And by the discriminant identity of the theorem \ref{thm:xi-terminal-zm-leyang-elliptic-structure}
\(
\mathrm{Disc}(P_{\mathrm{LY}})=-3^{9}\cdot 31^{2}\cdot 37
\)，
It can be seen that divergence value collisions may occur only at \(p\in\{3,31,37\}\); and the triple and double root cases can be distinguished by direct modular reduction.
On the other hand, the minimal model discriminant of \(E\) is \(\Delta(E)=37\) (Conclusion \ref{con:xi-terminal-zm-elliptic-normalization}), so \(p=37\) is the only bad reduced prime number;
And when \(p=2\) the transformation \(w=2y+1-2\lambda^2\) is no longer reversible and introduces an inseparable degeneration of the characteristic \(2\).
After removing the above finite prime set, the divergence index is relatively prime with \(p\) and the divergence values ​​are different from each other, thus covering the tame divergence and having good reduction. Certification completed.
\end{proof}

\begin{conclusion}[Fiber statistics implementation of modular form Hecke eigenvalue: total deviation identity]\label{con:xi-terminal-zm-hecke-eigenvalue-fiber-deviation}
For prime numbers \(p\notin\{2,37\}\), follow the conclusion \ref{con:xi-terminal-zm-fiber-root-counting-hecke}'s fiber root number \(N_p(y)\), and define the fiber difference function
\[
\delta_p(y):=N_p(y)-1.
\]
Then there is a strict identity
\[
\boxed{\;\sum_{y\in\FF_p}\delta_p(y)=-a_p\;}
\]
Where \(a_p=p+1-\#E(\FF_p)\) is the Frobenius trace of the elliptic curve \(E\).
Under the modular anchoring of the annotation \ref{rem:xi-terminal-elliptic-37a1-modular-anchor}, \(a_p\) is simultaneously equal to the new form of \(T_p\)-Hecke eigenvalue of the corresponding weight \(2\), level \(37\);
Therefore, this eigenvalue can be interpreted as the total deviation of "the number of fiber roots relative to the mean \(1\)" of the quartic spectral equation modulus \(p\).
\end{conclusion}
\begin{proof}[Proof sketch]
From the conclusion \ref{con:xi-terminal-zm-fiber-root-counting-hecke} the identity
\(
a_p=p-\sum_{y\in\FF_p}N_p(y)
\)，
Move the item and get it
\(
\sum_{y\in\FF_p}(N_p(y)-1)=-a_p
\)。
The Hecke eigenvalue interpretation is given by modular anchoring (see note \ref{rem:xi-terminal-elliptic-37a1-modular-anchor}). Certification completed.
\end{proof}

\begin{conclusion}[Specialization criterion falling into \texorpdfstring{$A_4$}{A4}: discriminant square condition and rational point constraint belonging to \texorpdfstring{$2$}{2} curve]\label{con:xi-terminal-zm-specialization-a4-square-criterion}
Take \(y_{0}\in\QQ\setminus\bigl(\{0,1\}\cup\{P_{\mathrm{LY}}(y)=0\}\bigr)\), and assume that the specialization four times \(\Pi(\lambda,y_{0})\in\QQ[\lambda]\) is irreducible on \(\QQ\).
then the criterion of judgment is established
\[
\boxed{\;
\mathrm{Gal}\bigl(\Pi(\lambda,y_{0})/\QQ\bigr)\subseteq A_{4}
\quad\Longleftrightarrow\quad
-y_{0}(y_{0}-1)P_{\mathrm{LY}}(y_{0})\in\QQ^{\times 2}\;}
\]
Equivalently, there exists \(w_{0}\in\QQ\) such that
\[
w_{0}^{2}=-y_{0}(y_{0}-1)P_{\mathrm{LY}}(y_{0}),
\]
That is \((y_{0},w_{0})\in C(\QQ)\), which belongs to \(2\) curve
\[
C:\quad w^{2}=-y(y-1)P_{\mathrm{LY}}(y)
\]
is the discriminant quadratic ridge curve of the conclusion \ref{con:xi-terminal-zm-discriminant-leyang}.
Therefore, all rational specialization parameters "falling into the intersection group" are strictly embedded in the set of rational points belonging to a \(2\) curve, and are therefore subject to strong rigid constraints in the arithmetical sense (for example, it can be deduced from the Faltings theorem that its set of rational points is finite).
\end{conclusion}
\begin{proof}[Proof sketch]
Under the assumption that \(y_{0}\) avoids divergent sets, the four roots of \(\Pi(\lambda,y_{0})\) are mutually different, so its specialized Galois group can be regarded as a transitive subgroup of \(S_{4}\).
For irreducible quartic, the Galois group becomes a chiasm group if and only if the discriminant is a rational square.
And from the conclusion \ref{con:xi-terminal-zm-discriminant-leyang},
\(
\mathrm{Disc}_{\lambda}\bigl(\Pi(\cdot,y_{0})\bigr)=-y_{0}(y_{0}-1)P_{\mathrm{LY}}(y_{0})
\)，
This leads to the equivalence shown; the characterization of the curve \(C\) is a straightforward rewriting. Certification completed.
\end{proof}



\begin{conclusion}[Lee--Yang divergence image characterization of cubic polynomials: Certificate of the minimum integer coefficient given by the resultant formula]\label{con:xi-terminal-zm-leyang-cubic-branch-image}
Bundle \(\Pi(\lambda,y)=0\) is deemed to be about \(\lambda\) and let the projection
\(
\pi:\mathcal C\to\mathbb A^1_y
\)
Depend on \((\lambda,y)\mapsto y\) is given.
but \(\pi\), the set of divergence values ​​is in \(\ZZ[y]\) is accurately characterized by the resultant expression:
\[
\boxed{\;
\mathrm{Res}_{\lambda}\!\bigl(\Pi(\lambda,y),\partial_{\lambda}\Pi(\lambda,y)\bigr)
=-y(y-1)\bigl(256y^3+411y^2+165y+32\bigr)\; }.
\]
Therefore, except for the trivial divergence value \(y\in\{0,1\}\), all non-trivial divergence values ​​are exactly cubic polynomials
\(
256y^3+411y^2+165y+32
=0
\)
of three roots; the only negative real root is the conclusion \ref{con:xi-terminal-zm-discriminant-leyang} of \(y_{\mathrm{LY}}\)。
The cubic factor is therefore not an empirical fitting object, but a projection \(\pi\) is like a minimum elimination certificate at the level of integer coefficients.
\end{conclusion}
\begin{proof}[Points of proof]
A divergence point exists if and only if \(\lambda\) makes \(\Pi(\lambda,y)=\partial_\lambda\Pi(\lambda,y)=0\)。
right \(\lambda\) can be decomposed by elimination; because \(\Pi\) first, the resulting formula is the same as \(\mathrm{Disc}_{\lambda}(\Pi)\) is consistent with the conclusion \ref{con:xi-terminal-zm-discriminant-leyang} The discriminant factorization of is completely consistent. Certification completed.
\end{proof}

\begin{conclusion}[point of divergence \texorpdfstring{$\lambda$}{lambda}-Identification of parameterization and Lee--Yang real roots]\label{con:xi-terminal-zm-leyang-lambda-cubic}
set up \((\lambda,y)\) satisfies the spectral curve \(\Pi(\lambda,y)=0\)。
Then all finite divergence points (that is, there are two-fold roots \(y\)) by simultaneous equations
\[
\Pi(\lambda,y)=0,\qquad \partial_{\lambda}\Pi(\lambda,y)=0
\]
Described by
\[
\boxed{\;(\lambda-1)(\lambda+1)\bigl(16\lambda^3-9\lambda^2+1\bigr)=0\;}
\]
And the corresponding divergence value is given by the explicit rational expression
\[
\boxed{\;
y=\frac{4\lambda^3-3\lambda^2-2\lambda+1}{4\lambda}
=\frac{(\lambda-1)(4\lambda^2+\lambda-1)}{4\lambda}\; }.
\]
in \(\lambda=1\) and \(\lambda=-1\) respectively correspond to the divergence values \(y=0\) and \(y=1\)；
The corresponding divergence values ​​of the remaining three roots are exactly the cubic equation \(256y^3+411y^2+165y+32=0\) of the three roots (conclusion \ref{con:xi-terminal-zm-leyang-cubic-branch-image}）。
In particular, the conclusion \ref{con:xi-terminal-zm-discriminant-leyang} The only negative real divergence value of \(y_{\mathrm{LY}}\) corresponds to the cubic equation \(16\lambda^3-9\lambda^2+1=0\) the only real root of \(\lambda_{\mathrm{LY}}\), and satisfy
%
thereby obtaining the \(y_{\mathrm{LY}}\) Equivalent characterization:
\[
\lambda_{\mathrm{LY}}\approx -0.2734348106524542,
\qquad y_{\mathrm{LY}}=y(\lambda_{\mathrm{LY}}),
\qquad 16\lambda_{\mathrm{LY}}^3-9\lambda_{\mathrm{LY}}^2+1=0.
\]
therefore \(y_{\mathrm{LY}}\) can be equivalently described as a rational mapping \(y=y(\lambda)\) in the third certificate \(16\lambda^3-9\lambda^2+1=0\) is the value at the unique real root.
\end{conclusion}
\begin{proof}[Points of proof]
Depend on \(\partial_{\lambda}\Pi(\lambda,y)=0\) can be solved
\(
y=(4\lambda^3-3\lambda^2-2\lambda+1)/(4\lambda)
\)（\(\lambda\neq 0\)）。
Substitute \(\Pi(\lambda,y)=0\) and clear the denominator to get
\(
(\lambda-1)(\lambda+1)(16\lambda^3-9\lambda^2+1)=0
\)；
in \(\lambda=\pm 1\) respectively correspond to the divergence values \(y=0,1\)；
The remaining three roots are characterized by cubic factors and \(y=y(\lambda)\) mapping falls into the conclusion \ref{con:xi-terminal-zm-leyang-cubic-branch-image} The cubic divergence image of , thus obtaining the unification of the Lee--Yang real divergence value and the accompanying spectral root. Certification completed.
\end{proof}

\begin{conclusion}[Lee--Yang Discriminant and quadratic subfield of cubic field:\texorpdfstring{$\QQ(\sqrt{-111})$}{Q(sqrt(-111))}]\label{con:xi-terminal-zm-leyang-cubic-arithmetic}
make
\(
g(y):=256y^3+411y^2+165y+32
\)。
Then its discriminant satisfies strict decomposition
\[
\boxed{\;\mathrm{Disc}(g)=-3^{9}\cdot 31^{2}\cdot 37\;}
\]
and \(g\) exist \(\QQ[y]\) is irreducible, so its Galois group is \(S_3\)。
remember \(L\) for \(g\), then its only quadratic subdomain is
\[
\boxed{\ \QQ(\sqrt{\mathrm{Disc}(g)})=\QQ(\sqrt{-111})\ },
\]
in \(-111\) for \(\mathrm{Disc}(g)\) is also the basic discriminant of the virtual quadratic field; its class number is
\[
h(-111)=8.
\]
It can be seen that the Lee--Yang three-dimensional arithmetic skeleton carries prime factors at the square free level at the same time \(37\) and \(3\), and pass \(\mathrm{Disc}(g)\) is uniquely locked in the split domain.
\end{conclusion}
\begin{proof}[Points of proof]
The discriminant factorization is calculated directly from the standard discriminant formula for cubic polynomials and is \(\ZZ\) is obtained by decomposing.
Can be excluded based on reasonable grounds \(g\), so that \(g\) is irreducible; for irreducible cubic, the discriminant is non-square if and only if the Galois group is \(S_3\)。
The secondary subdomain consists of \(\QQ(\sqrt{\mathrm{Disc}(g)})\) given; due to \(\mathrm{Disc}(g)\) The squared free part of \(-111\), so the quadratic domain is \(\QQ(\sqrt{-111})\)。
Number of classes \(h(-111)\) can be obtained by enumerating the reduced primitive positive definite binary quadratic form.
Reproducible verification script:\texttt{scripts/exp\_fold\_zm\_elliptic\_leyang\_arithmetic\_audit.py}. Certification completed.
\end{proof}

\begin{conclusion}[Lee--Yang three times in \texorpdfstring{$\FF_p$}{Fp} Decomposed density law on: by quadratic characteristics \texorpdfstring{$\left(\frac{-111}{p}\right)$}{(-111/p)} Complete control of the “just one” situation]\label{con:xi-terminal-zm-leyang-cubic-modp-splitting-density}
make
\[
g(y):=256y^3+411y^2+165y+32\in\ZZ[y],
\]
and take prime numbers \(p\nmid 3\cdot 31\cdot 37\)。
remember \(\varepsilon_p:=\left(\frac{-111}{p}\right)\in\{\pm1\}\) is a quadratic feature (Legendre symbol), and \(g\) module \(p\) turns into \(g_p\in\FF_p[y]\)。
but \(g_p\) satisfies:
\begin{itemize}
  \item if \(\varepsilon_p=-1\),but \(g_p\) must be \((1)(2)\) type, thus in \(\FF_p\) has exactly one root;
  \item if \(\varepsilon_p=+1\),but \(g_p\) or full split (\((1)(1)(1)\) type, three roots), or irreducible (\((3)\) type, rootless).
\end{itemize}
And the natural densities of the above three types of prime number sets are respectively
\[
\delta\bigl((1)(2)\bigr)=\frac12,\qquad
\delta\bigl((1)(1)(1)\bigr)=\frac16,\qquad
\delta\bigl((3)\bigr)=\frac13.
\]
therefore \(g_p\) exist \(\FF_p\) has at least one root (that is, the decomposition form is \((1)(2)\) or \((1)(1)(1)\)) The density of the set of prime numbers is
\[
\delta\bigl(g_p\ \text{exist}\ \FF_p\ \text{There are roots}\bigr)=\frac12+\frac16=\frac23.
\]
Discriminant quadratic ridge curve
\(
C_p:\ w^{2}=-y(y-1)g_p(y)
\)
In terms of , this is equivalent to the set of its divergence points at \(\FF_p\) divide by \(\{y=0,1,\infty\}\), additional \(\FF_p\)-Weierstrass point.
\end{conclusion}
\begin{proof}[Points of proof]
By conclusion \ref{con:xi-terminal-zm-leyang-cubic-arithmetic}，\(g\) exist \(\QQ\) and the Galois group of its split domain is \(S_3\), and has a unique quadratic resolution subdomain \(\QQ(\sqrt{-111})\)。
right \(p\nmid \mathrm{Disc}(g)\), Frobenius element \(\mathrm{Frob}_p\in S_3\) conjugate class with \(g_p\) has a one-to-one correspondence:
The identity class corresponds to total division \((1)(1)(1)\), transposition type correspondence \((1)(2)\), while the three-cycle class corresponds to irreducible \((3)\)。
On the other hand, symbolic homomorphism \(\mathrm{sgn}:S_3\to\{\pm1\}\)’s fixed domain is exactly the above-mentioned quadratic subdomain;
thus \(\mathrm{sgn}(\mathrm{Frob}_p)=\left(\frac{-111}{p}\right)=\varepsilon_p\)。
then \(\varepsilon_p=-1\) mandatory \(\mathrm{Frob}_p\) falls into the transposition conjugate class, thus \(g_p\) must be \((1)(2)\) type; while \(\varepsilon_p=+1\) is equivalent to \(\mathrm{Frob}_p\) is an even permutation, so that only identities or three cycles are possible, giving \((1)(1)(1)\) or \((3)\) two types.
The density is given by Chebotarev's theorem and is the size of the corresponding conjugate class divided by \(|S_3|=6\), that is, respectively \(3/6,1/6,2/6\)。
Reproducible verification script:\texttt{scripts/exp\_fold\_zm\_elliptic\_leyang\_cubic\_modp\_splitting\_audit.py}. Certification completed.
\end{proof}

\begin{conclusion}[prime number \texorpdfstring{$37$}{37} Free co-occurrence of squares: Ellipse discriminant, Lee--Yang bifurcation and second-order collision adjoint domain]\label{con:xi-terminal-37-squarefree-coupling}
in conclusion \ref{con:xi-terminal-zm-elliptic-normalization} Give the discriminant of the elliptic curve of the weight spectrum \(\Delta(E)=37\)；
in conclusion \ref{con:xi-terminal-zm-leyang-cubic-arithmetic} Give Lee--Yang three times \(g\) contains prime factors \(37\)。
On the other hand, let
\[
P_2(\lambda):=\lambda^3-2\lambda^2-2\lambda+2
\]
is the sum of second-order collision powers \(S_2(m)=\sum_{x\in X_m}d_m(x)^2\) of the minimal recursive characteristic polynomial (equivalent to \(y=1\) second-order collision kernel polynomial under the projection), then its discriminant is
\[
\boxed{\;\mathrm{Disc}\bigl(P_2\bigr)=148=4\cdot 37\;}
\]
Therefore, its accompanying quadratic domain is \(\QQ(\sqrt{37})\)。
on the other hand,\(E\) of \(2\)-twist point is given by \(2\)-division three times
\(
4X^{3}-4X+1=0
\)
Characterized, its discriminant is \(16\cdot 37\), so the only quadratic subdomain of its split domain is also \(\QQ(\sqrt{37})\)。
Therefore, at the level of square freedom, the same prime number \(37\) simultaneously appears in the weight spectrum ellipse discriminant, the Lee--Yang divergence cubic discriminant prime factor, and the discriminant adjoint domain of the second-order collision kernel, giving an auditable arithmetic coupling interface across three types of objects: "weight spectrum-divergence-collision".
\par
Furthermore, for the four discriminant objects that naturally appear in this article (\(2\)-division, divergence three times, Lee--Yang three times and Latt\`es critical six times) gives a unified enumeration of the free parts of the same odd prime square:
\[
\mathrm{Disc}\bigl(4X^3-4X+1\bigr)=2^{4}\cdot 37,
\qquad
\mathrm{Disc}\bigl(16X^3-9X^2+1\bigr)=-2^{2}\cdot 3^{3}\cdot 37,
\]
\[
\mathrm{Disc}\bigl(P_{\mathrm{LY}}(y)\bigr)=-3^{9}\cdot 31^{2}\cdot 37,
\qquad
\mathrm{Disc}\bigl(C(X)\bigr)=-2^{8}\cdot 37^{5},
\]
in \(P_{\mathrm{LY}}(y)=256y^3+411y^2+165y+32\),and \(C(X)=2X^{6}-10X^{4}+10X^{3}-10X^{2}+2X+1\) is the conclusion \ref{con:xi-terminal-zm-elliptic-lattes-doubling} The critical six times.
Therefore prime numbers \(37\) is an elliptic curve \(E\) is the only odd-prime bad reduction point, and it is also the source of odd-prime divergence that commonly appears in the above-mentioned divergence/critical generation domains.
\end{conclusion}
\begin{proof}[Points of proof]
\(\mathrm{Disc}(P_2)\) can be directly calculated from the cubic discriminant formula \(148\), its square free part is \(37\)。
The rest of the conclusions are drawn from the conclusion \ref{con:xi-terminal-zm-elliptic-normalization} and conclusion \ref{con:xi-terminal-zm-leyang-cubic-arithmetic} The discriminant calculations can be directly combined. Certification completed.
\end{proof}

% (User input window)

\begin{conclusion}[二阶碰撞三次与 \texorpdfstring{$E[2]$}{E[2]} 的三次坐标域：Tschirnhaus 生成元变换的显式同构]\label{con:xi-terminal-zm-collision-s2-e2-field-isomorphism}
沿用结论 \ref{con:xi-terminal-37-squarefree-coupling} 的二阶碰撞三次，并将自变量记为 \(\Lambda\)：
\[
P_2(\Lambda):=\Lambda^3-2\Lambda^2-2\Lambda+2\in\ZZ[\Lambda].
\]
另一方面，在结论 \ref{con:xi-terminal-zm-elliptic-normalization} 的短 Weierstrass 模型
\(
E:\ Y^2=X^3-X+\tfrac14
\)
上，非平凡 \(2\)-挠点的 \(X\)-坐标由 \(2\)-division 三次
\[
D_2(X):=4X^3-4X+1\in\ZZ[X]
\]
刻画。
取 \(\theta\) 为 \(D_2(X)=0\) 的任一根，并定义 Tschirnhaus 变换
\[
\boxed{\;\Lambda:=-2+2\theta+4\theta^2\;}
\]
则 \(P_2(\Lambda)=0\) 在 \(\QQ(\theta)\) 中成立，从而诱导域同构
\[
\boxed{\ \QQ(\Lambda)\cong \QQ(\theta)=:\QQ(E[2])\ }.
\]
因此二阶碰撞谱的三次代数数域与重量谱椭圆曲线的 \(2\)-挠坐标域在算术上被严格识别为同一对象；
并且
\(
\mathrm{Disc}(P_2)=4\cdot 37
\)
与
\(
\mathrm{Disc}(D_2)=2^{4}\cdot 37
\)
之平方自由部分一致性可由该显式同构统一解释。
\end{conclusion}
\begin{proof}[证明要点]
将 \(\Lambda=-2+2X+4X^2\) 代入 \(P_2(\Lambda)\) 并在 \(\QQ[X]\) 中对 \(D_2(X)=4X^3-4X+1\) 作带余除法，可得余式恒为零，从而当 \(D_2(\theta)=0\) 时必有 \(P_2(-2+2\theta+4\theta^2)=0\)。
等价地，结果式给出严格消元证书
\[
\mathrm{Res}_X\!\bigl(D_2(X),\ \Lambda-(-2+2X+4X^2)\bigr)=16\,P_2(\Lambda),
\]
因此 \(\Lambda\) 在 \(\QQ(\theta)\) 中满足同一三次极小多项式，诱导出所述三次域同构。
可复现核验脚本：\texttt{scripts/exp\_fold\_zm\_collision\_s2\_e2\_field\_isomorphism\_audit.py}。证毕。
\end{proof}

\begin{conclusion}[二阶碰撞分裂型与 \texorpdfstring{$E[2]$}{E[2]} 的 Frobenius 作用：模 \texorpdfstring{$p$}{p} 的可检验等价]\label{con:xi-terminal-zm-collision-s2-splitting-e2-torsion}
沿用结论 \ref{con:xi-terminal-zm-collision-s2-e2-field-isomorphism} 的记号，并取素数 \(p\neq 2,37\)。
则以下条件等价：
\begin{enumerate}
  \item 二阶碰撞三次 \(P_2(\Lambda)\) 在 \(\FF_p[\Lambda]\) 上完全分裂；
  \item \(2\)-division 三次 \(D_2(X)=4X^3-4X+1\) 在 \(\FF_p[X]\) 上完全分裂；
  \item 椭圆曲线 \(E/\FF_p\) 的三条非平凡 \(2\)-挠点全在 \(\FF_p\) 上有理，即
  \[
  E(\FF_p)[2]\cong (\ZZ/2\ZZ)^2.
  \]
\end{enumerate}
更进一步，\(P_2\)（等价地 \(D_2\)）在 \(\FF_p\) 上的分解型与 Frobenius 元在
\(
S_3\cong \mathrm{Aut}\bigl(E[2]\setminus\{O\}\bigr)
\)
中的共轭类一致，从而二阶碰撞谱的模素分裂行为可被精确译码为 \(E[2]\) 上的伽罗瓦表示型。
\par
与此相对，结论 \ref{con:xi-terminal-zm-leyang-ramification-critical-values} 中由三次式 \(16X^3-9X^2+1\) 所刻画的 Lee--Yang 非有理分歧轨道属于另一 \(S_3\)-型扩张，其二次预解域为 \(\QQ(\sqrt{-111})\)（见结论 \ref{con:xi-terminal-zm-leyang-cubic-arithmetic} 与结论 \ref{con:xi-terminal-zm-leyang-cubic-field-isomorphism}），一般并不与 \(E[2]\) 的分裂机制合并。
\end{conclusion}
\begin{proof}[证明要点]
由结论 \ref{con:xi-terminal-zm-collision-s2-e2-field-isomorphism} 的 Tschirnhaus 同构，可知 \(P_2\) 与 \(D_2\) 在 \(\QQ\) 上生成同一三次数域；
对 \(p\nmid 2\cdot 37\) 的素数，该数域在 \(p\) 处无分歧，因而 \(p\) 的分解型可由任一整系数生成多项式在 \(\FF_p\) 上的分解型读取，得到 \((1)\Longleftrightarrow(2)\)。
另一方面，在特征不为 \(2\) 的短 Weierstrass 形式下，非平凡 \(2\)-挠点恰对应 \(Y=0\)，其 \(X\)-坐标即为 \(D_2(X)=0\) 的三根；
故 \(D_2\) 在 \(\FF_p\) 上完全分裂当且仅当三条非平凡 \(2\)-挠点全为 \(\FF_p\)-点，即 \((2)\Longleftrightarrow(3)\)。
最后，Frobenius 在三条非平凡 \(2\)-挠点集合上的置换型与三次扩张的分解型之间的对应为标准事实。
可复现核验脚本：\texttt{scripts/exp\_fold\_zm\_collision\_s2\_e2\_field\_isomorphism\_audit.py}。证毕。
\end{proof}



